\section{候选对比方案与筛选结果}\label{app:candidate-schemes}

\subsection{与 JIT 时机和解释器路径相关的候选方向}

本附录集中归档第 5 章中未进入正文主线的候选方向。它们并不被称为“失败方案”，原因在于这些方向大多拥有
明确的实现对象和局部热点依据，只是尚未同时满足“命中正式 workload 主路径、形成 benchmark 级收益、具备稳定
解释闭环”这三个正文准入条件。因此，本文统一将其表述为候选方向、未保留方向或对比方案。

\begin{figure}[htbp]
    \centering
    \includegraphics[width=0.88\textwidth]{mybenchmark/failed-schemes-heatmap.pdf}
    \caption{附录候选方向在固定 baseline 下的归一化对比}
    \label{fig:failed-schemes-heatmap}
\end{figure}

图~\ref{fig:failed-schemes-heatmap} 给出了固定 baseline 下若干候选方向的统一比较结果。图中 \texttt{1.0}
表示 baseline，低于 \texttt{1.0} 说明该候选方向未跨过 benchmark 级收益门槛。需要强调的是，这张图的作用
并不是“展示所有负结果”，而是说明本文的正式方案来自系统性筛选，而非只挑选了几个成功样本。

\begin{table}[htbp]
    \centering
    \small
    \caption{JIT 时机与解释器路径相关候选方向的归档结论}
    \label{tab:appendix-jit-candidates}
    \begin{tabular}{>{\raggedright\arraybackslash}p{3.8cm}>{\raggedright\arraybackslash}p{2.2cm}>{\raggedright\arraybackslash}p{5.0cm}}
        \hline
        \textbf{候选方向} & \textbf{典型结果} & \textbf{未进入正文主线的原因} \\
        \hline
        \texttt{optimize\_threshold} 调整 & \texttt{pmd -11.23\%} & 作用于第二层 optimizing recompilation，收益窗口过晚 \\
        JIT OSR 热路径内联调整 & \texttt{pmd -17.78\%} & 假设更接近解释器/OSR 场景，无法解释正式 JIT 口径主瓶颈 \\
        Nterp 热点循环 OSR 重检查调整 & \texttt{mpegaudio -1.84\%} & 增加固定成本，未换来稳定收益 \\
        解释器已解析类型快路径 & \texttt{pmd -18.39\%} & 命中局部解析成本，但附加判断放大主路径开销 \\
        JIT debug-info 门控/延后 & 多数 workload 低于 baseline & 热点虽可见，但真实生效点与 benchmark 级收益缺少稳定闭环 \\
        \hline
    \end{tabular}
\end{table}

表~\ref{tab:appendix-jit-candidates} 说明，JIT 时机与解释器路径类候选方向大多存在同一问题：它们要么命中的是
第二层或更局部的运行路径，要么会通过额外检查、代码体积或时机扰动放大固定成本，最终难以形成与
\texttt{pmd} 这类短生命周期 workload 相匹配的正式结论。也正因如此，正文最终只保留了第一层
\texttt{warmup\_threshold} 微调这一最聚焦的控制流案例。

\subsection{与运行时热路径相关的候选方向}

\begin{table}[htbp]
    \centering
    \small
    \caption{运行时热路径相关候选方向的归档结论}
    \label{tab:appendix-runtime-candidates}
    \begin{tabular}{>{\raggedright\arraybackslash}p{4.0cm}>{\raggedright\arraybackslash}p{2.2cm}>{\raggedright\arraybackslash}p{4.8cm}}
        \hline
        \textbf{候选方向} & \textbf{典型结果} & \textbf{未进入正文主线的原因} \\
        \hline
        原始数组批量搬运展开 & \texttt{pmd -9.79\%} & 泛化复制路径不是正式 workload 的主导成本 \\
        \texttt{memcpy} 快速入口接线 & \texttt{pmd -11.95\%} & 目标 workload 不稳定落到该低层入口 \\
        CRC32 JNI 快路径 & \texttt{mpegaudio -14.28\%} & 目标热点占比仅约 \texttt{0.21\%}，优化对象选错 \\
        \texttt{DexFileVerifier} 连续栈容器替换 & \texttt{pmd -11.91\%} & 单一容器替换未覆盖完整启动期主链，故改写为正文阶段链案例 \\
        命令行参数/JIT init 相关方向 & 多数结果低于 baseline & 只压低局部 parser 开销，未命中主导成本 \\
        \hline
    \end{tabular}
\end{table}

与运行时热路径相关的候选方向进一步说明，局部热点的可见性并不等于正文价值。以 CRC32 JNI 快路径为例，
虽然它拥有明确函数入口，但目标 workload 上的占比过低；再如 \texttt{memcpy} 接线与原始数组批量搬运，
虽然都围绕“复制”展开，却没有像正文方案二那样持续命中字符串替换、拼接与 UTF-16 构造这些正式 workload 的
核心热路径。

值得注意的是，\texttt{DexFileVerifier} 单一容器替换方向在这一阶段也未能形成正式收益，因此并未以“单点修补”
的形式保留；相反，正文第 6.3 节采用的是更符合热点结构的阶段链组织，即沿 container probing、
UTF-8 校验与 LEB128 解码这条主链重写其叙事。

\subsection{未保留方向的共性原因}

综合图~\ref{fig:failed-schemes-heatmap} 与表~\ref{tab:appendix-jit-candidates}、
表~\ref{tab:appendix-runtime-candidates}，未保留方向主要呈现出三类共性原因。

\begin{enumerate}
    \item \textbf{只压低局部热点，但没有命中主导瓶颈。} 这类方向往往在 flamegraph 上可见，但热点占比不足，
    因而无法兑现为 benchmark 级收益。
    \item \textbf{假设前提不成立，或收益窗口过晚。} 第二层编译时机、OSR 调整等方向往往无法解决短生命周期
    workload 最关键的首层进入问题。
    \item \textbf{降低一种成本，却破坏更大的布局平衡。} 某些方向确实减少了某类局部开销，但同时引入了更大的
    代码体积、额外判断或复用损失，最终整体变慢。
\end{enumerate}

这些归档材料的意义，在于反向说明正文三项正式方案为何能够被保留下来：它们不仅命中了主路径，而且能够在固定
证据口径下稳定解释“为何有效、对谁有效、边界在哪里”。
%---------------------------------------------------------------------------%
